
\begin{abstract}
% Foreword
%     Context:
%     Need:
%     Task (what the authors did):
%     Object (what the document does):
% Summary
%     Findings:
%     Conclusion:
%     Perspectives

Este documento provee de los datos y la información para realizar el análisis de la deflexión de una viga de sección rectangular cargada al centro. Esto se hace mediante el calculo teórico de la deflexión de la viga y luego se realiza la comparación con los datos proporcionados originalmente. Luego se analizan los resultados mediante la diferencia relativa, esto evalúa qué tan parecidos son los valores que se midieron en la realidad respecto a los calculados con la fórmula teórica, para finalizar con una interpretación de estos y la conclusión respecto a estos.
 
\end{abstract}

\section{Introducción}
El análisis del comportamiento flexional de elementos estructurales continuos constituye un pilar fundamental en la ingeniería civil y mecánica. En el diseño y verificación de estructuras, la evaluación de las deformaciones transversales (deflexiones) bajo cargas de servicio es crítica para garantizar tanto los estados límites de servicio (funcionalidad, confort y estética) como la integridad estructural general.

El propósito de esta nota técnica es comparar y evaluar el comportamiento elástico-lineal de una viga simplemente apoyada sometida a un esquema de carga puntual centrada incremental, contrastando los desplazamientos verticales medidos experimentalmente frente a las formulaciones teóricas derivadas del modelo de Euler-Bernoulli. Para esto se entrega un conjunto de datos carga-deflexión para una viga simplemente apoyada con carga puntual centrada, junto con los parámetros geométricos y el módulo de elasticidad. 


\section{Desarrollo}
A continuación, se detalla el desarrollo metodológico y analítico para la caracterización mecánica del material de la viga a partir del ensayo de flexión en tres puntos.

%
\subsection{Objetivos del ensayo}
Caracterizar el comportamiento elástico del material sometido a esfuerzos de flexión directa mediante la aplicación de cargas puntuales incrementales.
%
\subsection{Caracterización de las Propiedades Mecánicas}
\begin{itemize}
\item Modulo de elasticidad: El módulo de elasticidad, también conocido como módulo de Young, es una propiedad física que mide la rigidez de un material y su capacidad para deformarse de forma elástica cuando se le aplica una fuerza.
\item Inercia: Es la magnitud física que mide la resistencia de un cuerpo u objeto a cambiar su estado de movimiento, ya sea al rotar o al trasladarse, o la resistencia de una sección estructural a deformarse por flexión.
\end{itemize}

%
\subsection{Definición de la geometría y formulación matemática}
Con el foco puesto en la caracterización del material y no en la complejidad geométrica, se adopta la geometría más sencilla posible que permita aislar la respuesta a flexión pura:

\begin{itemize}
    \item Configuración Geométrica: Viga prisma de sección rectangular constante, apoyada simplemente en sus extremos con una luz libre $L = 4.0\text{ m}$.
    \begin{figure}[hbtp]
    \centering
    \subfigure[]{\includegraphics[width=0.40\textwidth]{esquema_viga.png}}
    \caption{Esquema de la viga}
    \label{fig:fig_modelo}
    \end{figure}
    \item Dimensiones de la Sección:
    \begin{itemize}
        \item Ancho ($b$): $0.20\text{ m}$
        \item Altura ($h$): $0.40\text{ m}$
    \end{itemize}
    \item El momento de inercia es $I = \frac{b h^3}{12}$, donde $b$ es el ancho y $h$ la altura. 
    \item Deflexión teórica de Euler-Bernoulli en el centro, $\delta_{\text{máx}} = \frac{P L^3}{48 E I}$.
    Donde:
    \begin{itemize}
        \item P: Valor de la carga en $kN$
        \item L: Luz libre
        \item E: Módulo de elasticidad
        \item I: Inercia de la sección
    \end{itemize}
    
\end{itemize}
%

\subsection{Resultados y análisis}
En el Cuadro 1 se presentan los resultados de la deflexión teórica, comparando esta con la deflexión medida de acuerdo a cada carga aplicada.
\begin{table}[hbt]												
\caption{Comparación deflexión medida vs deflexión teórica}				
\centering	
\begin{tabular}{ccc} 											
\toprule														
\textit{Cargas [$kN$]} & \textit{Deflexión medida [$mm$]} & \textit{Deflexión teórica [$mm$]} \\
\midrule
0 & 0 & 0 \\
5 & 0,26 & 0,25 \\
10 & 0,5 & 0,5 \\
15 & 0,76 & 0,75 \\
20 & 1,01 & 1,00 \\
25 & 1,28 & 1,25 \\
30 & 1,54 & 1,50 \\
35 & 1,77 & 1,75 \\
40 & 2,05 & 2,00 \\								
\bottomrule
\end{tabular}
\end{table} 

A continuación se presenta el gráfico obtenido con los datos, comparando la deflexión medida y la deflexión teórica:
\begin{figure}[H]
    \centering
    \includegraphics[width=0.50\textwidth]{carga-deflexion.png}
    \caption{Gráfico Carga-Deflexión}
    \label{fig:fig_modelo2} % Sugerencia: cambia la etiqueta para no duplicarla con fig:fig_modelo
\end{figure}

Los resultados confirman que el comportamiento de la viga es estrictamente elástico lineal en el rango de $0$ a $40\text{ kN}$, manteniendo una flexibilidad constante de $\delta/P = 0,05\text{ mm/kN}$. La teoría de Euler-Bernoulli modela con elevada precisión el fenómeno real, presentando una diferencia relativa promedio del $1,9\%$ respecto a las mediciones experimentales. La ligera tendencia a registrar deflexiones medidas marginalmente mayores a las teóricas se atribuye a asentamientos menores en los apoyos y a la contribución no modelada de la deformación por cortante.


\section{Conclusiones}
A partir de la evaluación analítica y experimental realizada sobre la viga simplemente apoyada sometida a carga puntual centrada, se presentan las siguientes conclusiones principales:

\begin{itemize}
    \item Validación del Modelo Teórico: Se verificó una concordancia sobresaliente entre los valores de deflexión medidos y los calculados mediante la fórmula Euler-Bernoulli. La diferencia relativa promedio registrada fue de solo $1.88\%$, alcanzando una desviación máxima del $4.00\%$ (a los $5\text{ kN}$) y una diferencia absoluta no mayor a $0.05\text{ mm}$ en todo el rango de carga ($0\text{ a }40\text{ kN}$). Esto confirma que el modelo analítico simplificado resulta altamente preciso y adecuado para predecir la elástica del elemento en condiciones de servicio.
    \item Verificación de la Linealidad y Comportamiento Elástico: El material exhibió un comportamiento estrictamente elástico lineal a lo largo de todo el ensayo. La razón entre la deflexión medida y la carga aplicada ($\delta/P$) se mantuvo prácticamente constante en una pendiente experimental promedio de $0.051\text{ mm/kN}$ (frente a los $0.050\text{ mm/kN}$ teóricos), lo que garantiza la validez de la Ley de Hooke y demuestra que no se produjeron deformaciones plásticas ni fluencia del material bajo las cargas evaluadas.
\end{itemize}

\newpage
\section{Referencias}
Tablas de deflexión de vigas | MechaniCalc. (2026). Mechanicalc.Com. https://mechanicalc.com/reference/beam-deflection-tables

‌

